%%%
%
% $Autor: Wings $
% $Datum: 2021-05-14 $
% $Pfad: GitLab/MLEdgeComputer $
% $Dateiname: ESP32WiFi
% $Version: 4620 $
%
% !TeX spellcheck = de_DE/GB
% !TeX program = pdflatex
% !BIB program = biber/bibtex
% !TeX encoding = utf8
%
%%%

\chapter{Einführung: Wordclock}

\section{Allgemeines}
Die Messung der Zeit ist schon seit Jahrtausenden ein Thema für uns Menschen. Dabei hat sich die Art der Anzeige immer wieder verändert – von den ersten Sonnenuhren über mechanische Uhrwerke bis hin zu den heutigen digitalen Anzeigen. Eine besondere Form der Zeitanzeige ist die Wordclock.\\

\Mynote{Uhrentechnik}

Im Vergleich zu den klassischen analogen oder digitalen Uhren, die wir aus dem Alltag kennen und die meistens mit Zeigern oder Zahlen arbeiten \Mynote{uhren:technik Luebke, Anton, \textit{Die Uhr: Von der Sonnenuhr zur Atomuhr}, (Standardwerk zur Uhrentechnik/Geschichte).}
, zeigt eine Wordclock die Zeit in ganzen Worten an. Anstatt also Zahlen wie „10:15“ zu lesen, sieht man Sätze wie „Es ist viertel nach zehn“. Der große Unterschied ist hierbei die Lesbarkeit: Man muss nicht erst die Position der Zeiger interpretieren oder Zahlenwerte im Kopf umwandeln, sondern bekommt die Information direkt so, wie man sie auch aussprechen würde \Mynote{heisse}. Technisch wird das meistens so gelöst, dass hinter einer Platte mit Buchstaben eine Matrix aus LEDs sitzt, die genau die passenden Wörter beleuchtet \Mynote{heisse c't Magazin, \textit{Projekt: Word Clock}, \url{https://www.heise.de/select/ct/2016/24/1479984391777972}}.

\section{Applikations-Anforderungen}
Damit die Wordclock am Ende richtig funktioniert, muss die Hardware bestimmte Aufgaben erfüllen:
\begin{itemize}
	\item \textbf{Rechenleistung und Internet:} Als Herzstück wird ein Syde’s ESP32 Hypatia genutzt. Dieser Controller ist ideal, weil er schon WLAN eingebaut hat, was wir brauchen, um die Uhrzeit automatisch aus dem Internet zu ziehen.
	\item \textbf{Beleuchtung:} Für die Anzeige kommen WS2812B LED-Streifen zum Einsatz. Das sind digital adressierbare RGB-LEDs, bei denen wir jede einzelne LED im Streifen separat steuern können, was die Programmierung der Buchstaben deutlich einfacher macht.
	\item \textbf{Stromversorgung:} Da die vielen LEDs zusammen recht viel Strom verbrauchen können, nutzen wir einen Stepdown-Konverter. Dieser sorgt dafür, dass die Spannung stabil bleibt und die Bauteile nicht beschädigt werden.
\end{itemize}

\section{Herausforderungen}
Bei der Planung und dem Bau gibt es ein paar Punkte, die man im Blick behalten muss:
\begin{itemize}
	\item \textbf{Die richtige Uhrzeit:} Da wir ohne eine Real Time Clock arbeiten, muss die Uhrzeit über das WLAN per NTP-Protokoll regelmäßig abgefragt werden. 
	\item \textbf{Hitzeentwicklung:} Wenn viele LEDs gleichzeitig hell leuchten, entsteht Wärme. Das Gehäuse muss so gebaut sein, dass sich diese Hitze nicht anstaut.
	\item \textbf{Licht-Trennung:} Damit ein Buchstabe nicht den Nachbarbuchstaben mit beleuchtet, muss man im Gehäuse saubere Trennwände zwischen den LEDs einplanen.
	\item \textbf{Stromversorgung:} Da die Vielen LEDs zusammen eine dementsprechend große Leistung benötigen ist die Auslegung der Stromversorgung im Hinblick auf die Kompatibilität mit dem Hypatiaboard zu beachten.  
\end{itemize}

\section{Aufbau des Berichtes} \Mynote{muss ausfürlicher}
Dieser Bericht ist in mehrere Teile gegliedert, um den gesamten Weg vom Konzept bis zur fertigen Uhr zu zeigen:
Nach dieser \textbf{Einleitung (Part 1)} geht es in \textbf{Part 2} um die Tools, wie zum Beispiel die Arduino IDE. In \textbf{Part 3} wird die Hardware genau beschrieben – angefangen beim Prozessor bis hin zu den Sensoren und dem ersten Testcode. Die Software und der finale Code werden in \textbf{Part 4} erklärt. Danach folgt in \textbf{Part 5} ein Überblick über das Gehäuse und in \textbf{Part 6} die eigentliche Anwendung mit den Abläufen. Zum Schluss ziehen wir in \textbf{Part 7} ein Fazit und schauen in der Evaluation nach, ob wir alle Anforderungen und Herausforderungen aus der Einleitung erfolgreich gemeistert haben.



